\documentclass[runningheads]{llncs}
%
\usepackage[T1]{fontenc}
% T1 fonts will be used to generate the final print and online PDFs,
% so please use T1 fonts in your manuscript whenever possible.
% Other font encondings may result in incorrect characters.
%
\raggedbottom
\usepackage{graphicx}
% Used for displaying a sample figure. If possible, figure files should
% be included in EPS format.
\usepackage{float}
% If you use the hyperref package, please uncomment the following two lines
% to display URLs in blue roman font according to Springer's eBook style:
\usepackage{color}
\usepackage{hyperref}
\renewcommand\UrlFont{\color{blue}\rmfamily}
\urlstyle{rm}
%
\begin{document}
%
\title{Report for Assignment 1}
%
%\titlerunning{Abbreviated paper title}
% If the paper title is too long for the running head, you can set
% an abbreviated paper title here
%
\author{Mauritz Almgren\inst{1}\\ \email{manz24@student.bth.se}}

%
\authorrunning{M. Almgren}
% First names are abbreviated in the running head.
% If there are more than two authors, 'et al.' is used.
%
\institute{Blekinge Tekniska Högskola\\Valhallavägen 10, 371 79 Karlskrona, Sweden\\
\email{info@bth.se}}
%
\maketitle              % typeset the header of the contribution
%

\section{Introduction}

The purpose of this report is to analyse 2 types of agents, reflex- and utility-based, and to implement and compare search algorithms, Breadth First Search and Depth First Search based on their ability to solve a maze problem.

\section{Problem 1 Shopping Assistant Agents}

The problem that the agents have to solve is to help the customer buy items from a catalog based on the customers' budget and the items price and satisfaction. 

The agents that were used to solve this problem were reflex and utility based. 

A reflex agent is a very simple agent that decides what to do based on "reflexes" which can be regarded as if-then statements. The reflex agent that was implemented for this problem used these if-then statements:
\begin{verbatim}
    if amount + item.price <= budget:
            if item.price <= 5:
                then buy item
            elif item.satisfaction > 8:
                then buy item
\end{verbatim}

The implementation is as such that the agent checks for every items price if it is under or equals to five and if so buys it, or if the items satisfaction is over eight it also buys it. \newline

A utility agent is more sophisticated as compared to the reflex agent in its decision-making. It uses different sets of conditions to find the maximum of some given parameters for all possible combinations of the items in the basket. The utility agent that was implemented to solve the problem used these conditions:
\begin{verbatim}
    for r in combinations(items, i):
            comb_price = sum(item.price for item in r)
            comb_satisfaction = sum(item.satisfaction for item in r)
            if comb_price <= budget:
                combs.append(Combination(r, comb_price, comb_satisfaction))
    best_comb = max(combs, key=lambda comb: (comb.total_satisfaction, -comb.total_price))
\end{verbatim}

The key parameters in this utility agent is the total satisfaction for a given combination and it's negative total price. The total satisfaction gives us the combination of items with the highest total satisfaction. The negative total price gives the cheapest combination for the highest total satisfaction, this is because the negative gives us a way to find the lowest price for combinations with the same amount of total satisfaction.

\section{Problem 2 Maze Solving}

The problem that the search algorithms have to solve is to find the shortest path in a given maze. The maze that was given in this problem is as provided below:


\end{document}